%=====================================================================
%  II. METODOLOGI PENELITIAN
%=====================================================================
\section{METODOLOGI PENELITIAN}

\subsection{Jenis Penelitian}

Penelitian ini merupakan penelitian analitik-komputasional dengan pendekatan
deduktif-matematis. Formulasi fisika diturunkan secara analitik dari prinsip dasar
tanpa melibatkan pengambilan data eksperimental, kemudian dilanjutkan dengan
komputasi numerik sederhana untuk memvisualisasikan besaran yang diperoleh.
Pelaksanaan penelitian diawali dengan studi literatur mengenai teori relativitas
umum, karakteristik metrik Schwarzschild, dan perturbasi medan uji pada latar
belakang ruang-waktu yang melengkung.

\subsection{Objek dan Asumsi Penelitian}

Objek penelitian ini adalah medan skalar uji $\phi$ yang merambat pada latar
belakang ruang-waktu Schwarzschild. Medan skalar diperlakukan sebagai medan uji,
yakni amplitudo gangguannya dianggap cukup kecil sehingga tidak memberikan pengaruh
balik terhadap geometri latar belakang. Konsekuensinya, metrik latar belakang
diperlakukan tetap dan hanya persamaan gerak medan skalar yang diselesaikan.
Seluruh perhitungan menggunakan \textit{natural unit} $G = c = 1$, sehingga massa
lubang hitam $M$ memiliki dimensi yang sama dengan koordinat radial $r$. Konstanta
Kosmologi diambil bernilai nol, sehingga ruang-waktu bersifat datar secara
asimtotik.

\subsection{Tahapan Penelitian}

Penelitian dilaksanakan melalui lima tahap. Pertama, dituliskan persamaan gerak
medan skalar dalam ruang-waktu lengkung, yaitu persamaan Klein-Gordon, yang
diperoleh dari prinsip aksi minimal. Kedua, komponen tensor metrik Schwarzschild
disubstitusikan ke dalam persamaan tersebut sehingga diperoleh bentuk eksplisitnya
dalam koordinat $(t, r, \theta, \varphi)$. Ketiga, dilakukan separasi variabel
dengan memanfaatkan sifat statis dan simetri bola latar belakang, sehingga
diperoleh persamaan radial. Keempat, persamaan radial ditransformasikan ke dalam
koordinat \textit{tortoise} dan disubstitusi dengan fungsi radial yang telah
diskalakan, sehingga diperoleh persamaan master beserta bentuk eksplisit potensial
efektifnya. Kelima, potensial efektif yang diperoleh dievaluasi dan divisualisasikan
secara numerik menggunakan bahasa pemrograman Python dengan pustaka NumPy dan
Matplotlib, untuk variasi bilangan momentum sudut $\ell$, massa lubang hitam $M$,
dan massa medan skalar $\mu$.

\subsection{Persamaan yang Digunakan}

Ruang-waktu Schwarzschild dideskripsikan oleh elemen garis
\begin{equation}
  ds^{2} = -f(r)\,dt^{2} + \frac{dr^{2}}{f(r)}
           + r^{2}\left(d\theta^{2} + \sin^{2}\theta\, d\varphi^{2}\right),
  \label{eq:metrik}
\end{equation}
dengan fungsi metrik
\begin{equation}
  f(r) = 1 - \frac{2M}{r},
  \label{eq:fungsi-metrik}
\end{equation}
sehingga horizon peristiwa terletak pada $r_{h} = 2M$, yakni akar dari $f(r) = 0$.

Dinamika medan skalar uji tak bermassa pada latar belakang tersebut mematuhi
persamaan Klein-Gordon dalam ruang-waktu lengkung,
\begin{equation}
  \Box \phi = \frac{1}{\sqrt{-g}}\,
  \partial_{\mu}\left(\sqrt{-g}\, g^{\mu\nu} \partial_{\nu} \phi\right) = 0,
  \label{eq:kg}
\end{equation}
sedangkan untuk medan skalar bermassa ditambahkan suku massa $\mu$ sehingga
\begin{equation}
  \partial_{\mu}\left(\sqrt{-g}\, g^{\mu\nu} \partial_{\nu} \phi\right)
  - \sqrt{-g}\,\mu^{2}\phi = 0,
  \label{eq:kg-bermassa}
\end{equation}
dengan $\mu$ menyatakan massa medan skalar itu sendiri, yang perlu dibedakan dari
massa lubang hitam $M$.

Separasi variabel dilakukan dengan ansatz
\begin{equation}
  \phi(t, r, \theta, \varphi)
  = e^{-i\omega t}\, R(r)\, Y_{\ell m}(\theta, \varphi),
  \label{eq:ansatz}
\end{equation}
dengan $Y_{\ell m}$ adalah harmonik sferis dan $\ell$ bilangan momentum sudut.
Adapun koordinat \textit{tortoise} $r_{*}$ didefinisikan melalui hubungan
\begin{equation}
  \frac{dr_{*}}{dr} = \frac{1}{f(r)},
  \label{eq:tortoise-def}
\end{equation}
yang untuk latar belakang Schwarzschild memberikan
\begin{equation}
  r_{*} = r + 2M \ln\left(\frac{r}{2M} - 1\right).
  \label{eq:tortoise}
\end{equation}
